\scanpage{169}
\mspage{85}
% archive sheet number: 168
% lines 1-3
Rappelons que la cofinalité des foncteurs (4.4.6.3) \inserted{quel que soit le $j\in\operatorname{Ob}J$} implique le fait que $\varphi$ est une limite inductive du foncteur \inserted{composé}

% cancelled attempts above the live diagram
\deleted{\diagram{\begin{tikzcd}[ampersand replacement=\&]
I \arrow[r,"\unsure{\psi}"] \arrow[d,equal] \& \underline{\operatorname{Hom}}(J,C) \arrow[r,"\alpha^J"] \& \unsure{\underline{\operatorname{Hom}}(J,\operatorname{Ind}'_{\pi}(C))} \\
\underline{\operatorname{Hom}}(J,C)_{/\varphi} \& \underline{\operatorname{Hom}}(J,C) \arrow[r,"\alpha^J"] \& \unsure{\underline{\operatorname{Hom}}(J,\operatorname{Ind}'_{\pi}(C))}
\end{tikzcd}}}

% main diagram
\diagram{\begin{tikzcd}[ampersand replacement=\&,column sep=large]
I \arrow[r,"\unsure{\psi^J|\psi_J}"] \arrow[d,equal] \& \underline{\operatorname{Hom}}(J,C) \arrow[r,"\alpha^J"] \arrow[dd,dotted,no head,"\varepsilon_j^*"] \& \underline{\operatorname{Hom}}(J,\operatorname{Ind}'_{\pi}(C)) \arrow[dd,dotted,"\varepsilon_j^{*\prime}"] \\
\underline{\operatorname{Hom}}(J,C)_{/\varphi} \arrow[d,"\unsure{\varepsilon_j^{\varphi}|\xi_j^{\varphi}}"'] \arrow[dr,phantom,"\deleted{\unsure{\varphi_j}}"] \& \& \\
C_{/\varphi(j)} \arrow[r,dashed,"\psi_j"] \& C \arrow[r,dashed,"\alpha"] \& \operatorname{Ind}'_{\pi}(C)
\end{tikzcd}}
\[\tag*{(4.4.6.4)}\vphantom{I}\]

\inserted{qui est l'ens. ordonné filtrant $I$ \deleted{\unsure{(plus)}} \unsure{précisément}}

% large cancelled paragraph, lines 8-13
\deleted{$\varphi$ sera $=$ la limite inductive dans la catégorie plus grande $\underline{\operatorname{Hom}}(J,\widehat C\deleted{\unsure{(C)|C}})$, admettant $\underline{\operatorname{Hom}}(J,\operatorname{Ind}_{\pi}C)$ comme sous-catégorie pleine; car \unsure{cette dernière assertion signifie}}

% lines 14-22
\unsure{plus précisément} \unsure{est que|que} $\varphi$ est la $\varinjlim$ ``argument par argument $j$'', ce qui signifie que pour tout $j\in\operatorname{Ob}J$, $\varphi(j)=\unsure{\varepsilon_j^*|\varepsilon_j^{*\prime}}(\varphi)$ est limite inductive des $\varepsilon_j^{*\prime}(\alpha^J\psi)=(\alpha\unsure{\psi_j|\varepsilon_j^*})\varepsilon_j^{\varphi}$.

\deleted{$\unsure{\varphi}$} (Or on sait \inserted{en effet} que $\varphi(j)$ est la limite de $\alpha\psi$, puisque $\varphi(j)\in\operatorname{Ob}\widehat C$ est ind-représentable, donc aussi celle des $(\alpha\psi_j)\varepsilon_j^{\varphi}$ si $\varepsilon_j^{\varphi}$ est cofinal.)

% final lines continue on next scan
Ainsi, admettant la cofinalité \inserted{de (4.4.6.3)}, \deleted{\unsure{dans}} et introduisant $\psi_J$ de (4.4.6.3)
